\documentclass[../DoAn.tex]{subfiles}
\begin{document}
\label{ch:nen_tang}
% =============================================================
% CHƯƠNG 1 - NỀN TẢNG LÝ THUYẾT          Mục tiêu độ dài: ~10 trang
% =============================================================

Chương này hệ thống hoá các kiến thức nền tảng phục vụ cho phương pháp đề xuất: học tăng cường sâu, mạng nơ-ron đồ thị và cơ chế attention, ba khối nền tảng được sử dụng trong đồ án. Phần khảo sát các hướng tiếp cận hiện có cho bài toán và phát biểu hình thức của bài toán được trình bày ở Chương~\ref{ch:bai_toan}.

\section{Học tăng cường sâu}
\label{sec:rl}

\subsection{Quá trình quyết định Markov}

Học tăng cường được hình thức hoá qua \textit{Quá trình quyết định Markov} (Markov Decision Process - MDP), một bộ năm
\[
    \mathcal{M} = (\mathcal{S}, \mathcal{A}, \mathcal{P}, \mathcal{R}, \gamma),
\]
trong đó $\mathcal{S}$ là không gian trạng thái, $\mathcal{A}$ là không gian hành động, $\mathcal{P}(s' \mid s, a)$ là xác suất chuyển trạng thái, $\mathcal{R}(s, a)$ là hàm phần thưởng, và $\gamma \in [0, 1]$ là hệ số chiết khấu. Một \textit{chính sách} $\pi(a \mid s)$ là phân phối xác suất trên hành động cho mỗi trạng thái. Mục tiêu của RL là tìm chính sách $\pi^*$ tối đa hoá \textit{kỳ vọng lợi ích tích luỹ}\cite{sutton2018rl}:
\begin{equation}
    J(\pi) = \mathbb{E}_{\tau \sim \pi}\left[ \sum_{t=0}^{T} \gamma^t \mathcal{R}(s_t, a_t) \right],
\end{equation}
trong đó $\tau = (s_0, a_0, s_1, a_1, \ldots)$ là quỹ đạo tương tác. Việc kết hợp RL với mạng nơ-ron sâu (deep reinforcement learning) đã đạt nhiều thành tựu nổi bật, tiêu biểu là kiểm soát ở mức con người trên các trò chơi Atari\cite{mnih2015dqn}.

Trong bối cảnh VNE, một episode tương ứng với quá trình nhúng một VNR. Trạng thái mô tả mạng vật lý hiện tại và VNR đang xét, hành động là ánh xạ một nút ảo tới một nút vật lý hoặc xếp hạng các nút/liên kết ảo, và phần thưởng được cấp khi VNR đã được ánh xạ hoàn chỉnh và đo bằng tỉ số R/C.

\subsection[Gradient chính sách và phương pháp tác tử-phê bình]{Gradient chính sách và phương pháp tác tử-phê bình (Actor-Critic)}

\textit{Policy gradient} là họ thuật toán RL học chính sách trực tiếp bằng cách tham số hoá $\pi_\theta(a \mid s)$ với tham số $\theta$ và cập nhật theo gradient của $J(\theta)$. Định lý policy gradient\cite{sutton2018rl} chỉ ra:
\begin{equation}
    \nabla_\theta J(\theta) = \mathbb{E}_{\tau \sim \pi_\theta}\left[ \sum_t \nabla_\theta \log \pi_\theta(a_t \mid s_t) \cdot G_t \right],
\end{equation}
với $G_t = \sum_{k=t}^{T} \gamma^{k-t} \mathcal{R}(s_k, a_k)$. REINFORCE\cite{williams1992reinforce} là phiên bản đơn giản nhất, ước lượng kỳ vọng bằng trung bình mẫu Monte Carlo. Hạn chế chính của REINFORCE là phương sai gradient cao. Để giảm phương sai mà không gây bias, có thể trừ một \textit{baseline} $b(s)$ độc lập với hành động:
\begin{equation}
    \nabla_\theta J(\theta) = \mathbb{E}\left[ \sum_t \nabla_\theta \log \pi_\theta(a_t \mid s_t) \cdot (G_t - b(s_t)) \right].
\end{equation}
Phương án phổ biến nhất là dùng hàm giá trị $V_\phi(s)$ làm baseline, dẫn đến kiến trúc \textit{Actor-Critic} với hai mạng song song: \textit{actor} là chính sách $\pi_\theta(a\mid s)$ quyết định hành động, còn \textit{critic} là hàm giá trị $V_\phi(s)$ ước lượng chất lượng trạng thái và được học bằng cách hồi quy về lợi ích thực tế (mất mát bình phương). Lợi thế $A(s,a)=G_t - V_\phi(s)$ thay cho $G_t$ giúp giảm phương sai gradient. Đồ án này sử dụng baseline phê bình $V_\phi(s)$ học được (kiến trúc actor-critic) thay cho baseline trung bình đơn giản, nhằm giảm phương sai mạnh hơn trong các kịch bản phần thưởng cấp ở cuối episode.

\subsection{Tối ưu hoá chính sách tiệm cận (PPO)}

\textit{Proximal Policy Optimization} (PPO)\cite{schulman2017ppo} là một thuật toán actor-critic cải tiến, cho phép tái sử dụng dữ liệu rollout qua nhiều vòng gradient mà vẫn ổn định. PPO thay mục tiêu gradient chính sách bằng một \textit{hàm mục tiêu thay thế có cắt} (clipped surrogate) dựa trên tỉ số tầm quan trọng $\rho_t(\theta) = \pi_\theta(a_t\mid s_t)/\pi_{\theta_{\text{old}}}(a_t\mid s_t)$:
\begin{equation}
    L^{\text{CLIP}}(\theta) = \mathbb{E}_t\Big[ \min\big( \rho_t(\theta)\,\hat{A}_t,\; \mathrm{clip}(\rho_t(\theta), 1-\epsilon, 1+\epsilon)\,\hat{A}_t \big) \Big],
    \label{eq:ppo-clip}
\end{equation}
với $\epsilon$ là biên cắt (thường $0{,}2$). Phép cắt ngăn chính sách trôi quá xa khỏi chính sách cũ trong một bước cập nhật, đóng vai trò một ràng buộc vùng tin cậy mềm. Đồ án sử dụng kiến trúc actor-critic làm phương pháp tinh chỉnh chính, đồng thời cài đặt biến thể PPO có cắt để so sánh (Chương~\ref{ch:thuc_nghiem}).

\section{Mạng nơ-ron đồ thị}
\label{sec:gnn}

\subsection{Cơ chế truyền tin}

Mạng nơ-ron đồ thị (Graph Neural Network - GNN) học biểu diễn của các nút bằng cơ chế \textit{truyền tin} (message passing). Tại mỗi lớp $\ell$, vector biểu diễn $h_v^{(\ell)}$ của mỗi nút $v$ được cập nhật từ biểu diễn của các láng giềng:
\begin{equation}
    h_v^{(\ell+1)} = \text{UPDATE}^{(\ell)}\left( h_v^{(\ell)},\; \text{AGGREGATE}^{(\ell)}\left( \{h_u^{(\ell)} : u \in \mathcal{N}(v)\} \right) \right).
\end{equation}
Sau $L$ lớp truyền tin, mỗi nút ``thấy'' được thông tin từ các láng giềng trong bán kính $L$ chặng. Các biến thể GNN khác nhau chủ yếu ở hai hàm AGGREGATE và UPDATE.

\subsection{Graph Convolutional Network}

\textit{Graph Convolutional Network} (GCN)\cite{kipf2017gcn} sử dụng phép trung bình có trọng số trên láng giềng:
\begin{equation}
    H^{(\ell+1)} = \sigma\left( \tilde{D}^{-1/2} \tilde{A} \tilde{D}^{-1/2} H^{(\ell)} W^{(\ell)} \right),
    \label{eq:gcn}
\end{equation}
trong đó $\tilde{A} = A + I$ là ma trận kề có self-loop, $\tilde{D}$ là ma trận đường chéo độ bậc, $H^{(\ell)} \in \mathbb{R}^{|V| \times d_\ell}$ là ma trận biểu diễn, $W^{(\ell)}$ là tham số học được, và $\sigma$ là hàm kích hoạt phi tuyến (đồ án dùng ReLU). Đồ án sử dụng GCN hai lớp với $H^{(0)} = X$ là ma trận đặc trưng đầu vào của các nút vật lý. Ma trận kề $A$ được trọng số hoá theo tỉ lệ băng thông còn lại trên dung lượng ban đầu, cho phép mã hoá thông tin trạng thái tài nguyên động vào quá trình truyền tin.

\subsection{Graph Attention Network}

\textit{Graph Attention Network} (GAT)\cite{velickovic2018gat} thay phép trung bình đồng đều bằng phép tổ hợp có trọng số attention học được:
\begin{equation}
    h_v^{(\ell+1)} = \sigma\left( \sum_{u \in \mathcal{N}(v) \cup \{v\}} \alpha_{vu}^{(\ell)} \cdot W^{(\ell)} h_u^{(\ell)} \right),
\end{equation}
với hệ số attention
\begin{equation}
    \alpha_{vu}^{(\ell)} = \frac{\exp\left( \text{LeakyReLU}(\mathbf{a}^\top [W h_v \| W h_u]) \right)}{\sum_{u' \in \mathcal{N}(v) \cup \{v\}} \exp\left( \text{LeakyReLU}(\mathbf{a}^\top [W h_v \| W h_{u'}]) \right)}.
\end{equation}
GAT có khả năng biểu diễn cao hơn GCN nhờ trọng số học được, đặc biệt khi đồ thị không đồng nhất về vai trò các láng giềng. Tuy nhiên, GCN có ưu điểm về số tham số ít, tốc độ huấn luyện nhanh, và đủ biểu diễn cho các tô-pô mạng vật lý nhỏ - vừa trong VNE. Đồ án này lựa chọn GCN làm bộ mã hoá chính để giữ kiến trúc nhỏ gọn (toàn bộ chính sách chỉ khoảng $36{,}300$ tham số), đồng thời thảo luận lại các đánh đổi với GAT trong phân tích ablation ở Chương~\ref{ch:thuc_nghiem}.

\section{Cơ chế chú ý cho lựa chọn ứng viên}
\label{sec:attention}

Cơ chế \textit{scaled dot-product attention}\cite{vaswani2017attention} là khối toán học chuẩn cho việc ``một truy vấn chọn từ nhiều khoá''. Cho truy vấn $q \in \mathbb{R}^d$ và tập khoá $K = (k_1, \ldots, k_n)$, điểm số cho mỗi khoá là
\begin{equation}
    \text{score}(q, k_i) = \frac{q^\top k_i}{\sqrt{d}}.
\end{equation}
Hệ số $\sqrt{d}$ giúp ổn định gradient khi $d$ lớn. Trong bài toán VNE, mỗi nút ảo đóng vai trò truy vấn, mỗi nút vật lý ứng viên đóng vai trò khoá, và điểm số sau khi softmax biểu diễn xác suất chọn từng ứng viên. Đồ án mở rộng cơ chế này bằng cách kết hợp một nhánh MLP dư (residual) để học các tương tác phi tuyến phức tạp giữa đặc trưng của nút ảo và nút vật lý, và sử dụng \textit{feasibility mask} cứng để loại bỏ các ứng viên không thoả mãn ràng buộc CPU ngay tại bước tính điểm. Chi tiết kiến trúc được trình bày ở Chương~\ref{ch:de_xuat}.

\section{Tối ưu hoá bầy đàn (Particle Swarm Optimization)}
\label{sec:pso}

Tối ưu hoá bầy đàn (Particle Swarm Optimization - PSO)\cite{kennedy1995pso} là một thuật toán tối ưu hoá metaheuristic lấy cảm hứng từ hành vi bầy đàn trong tự nhiên. PSO duy trì một \textit{quần thể} gồm $N$ \textit{hạt} (particle), mỗi hạt biểu diễn một lời giải ứng viên bằng một vector vị trí $x_i$ trong không gian tìm kiếm và di chuyển với vận tốc $v_i$. Mỗi hạt ghi nhớ vị trí tốt nhất của riêng nó ($p_i$, personal best), còn cả quần thể chia sẻ vị trí tốt nhất toàn cục ($g$, global best). Tại mỗi vòng lặp, vận tốc và vị trí của hạt được cập nhật theo:
\begin{align}
    v_i &\leftarrow w\,v_i + c_1 r_1 (p_i - x_i) + c_2 r_2 (g - x_i), \label{eq:pso-vel}\\
    x_i &\leftarrow x_i + v_i, \label{eq:pso-pos}
\end{align}
trong đó $w$ là \textit{hệ số quán tính} (inertia weight)\cite{shi1998modified} điều tiết mức duy trì hướng di chuyển cũ, $c_1$ và $c_2$ lần lượt là \textit{hệ số nhận thức} (cognitive) và \textit{hệ số xã hội} (social) điều tiết lực hút về $p_i$ và $g$, còn $r_1, r_2 \sim \mathcal{U}(0,1)$ là các số ngẫu nhiên tạo tính khám phá. Chất lượng mỗi hạt được đánh giá qua một \textit{hàm thích nghi} (fitness function) tuỳ theo bài toán.

PSO phù hợp với các không gian tìm kiếm rời rạc, gồ ghề và không khả vi như bài toán ánh xạ nút trong VNE, nơi khó áp dụng các phương pháp dựa trên gradient. Trong đồ án, PSO được dùng làm cơ chế tìm kiếm lời giải ánh xạ nút trong heuristic MP-VNE, và đặc biệt trong bước \textit{triển khai suy diễn} của CARL-VNE: hàm thích nghi của PSO được thiên lệch theo xác suất của chính sách đã học, kết hợp với chiến lược đa khởi tạo (multi-restart) để khai thác đồng thời nhiều vùng triển vọng (chi tiết ở Chương~\ref{ch:de_xuat}).

% Kết chương
Chương này đã hệ thống hoá các nền tảng học tăng cường (MDP, gradient chính sách, actor-critic, PPO), các khối GNN, cơ chế attention và tối ưu hoá bầy đàn (PSO). Chương tiếp theo khảo sát các hướng tiếp cận hiện có và phát biểu hình thức bài toán nhúng mạng ảo đa miền.

\end{document}
